\section{Limitations}\label{sec:limitations}

\paragraph{What is proved.} The limit theorems concern the first tree of a boosted ensemble with exact search; histogram search is covered under an assumption on XGBoost's quantile sketch that we have not proved. \Cref{thm:scaling} restricts thresholds to a fixed window, and the behaviour on XGBoost's full window is \cref{conj:drift}. The level results (\cref{prop:level,prop:rounds}) hold for any rank-based learner and any number of trees, but only for strong nulls: a feature that is irrelevant yet correlated with a relevant one is not covered. The power of the shadow rule on a full ensemble is \cref{conj:ensemble}, which the experiments support cell by cell but which we have not proved. \Cref{thm:limit} gives no rate of convergence.

\paragraph{What the experiments cover.} The main sweeps use Gaussian location shifts with conditionally independent features; \cref{sec:robustness} departs from this in eight ways, and \cref{sec:realdata} uses five real tables in which only the appended null columns are known to be irrelevant, so the false negative rates there rest on a pseudo-truth. We study binary targets; a rare class in a multi-class problem reduces to the same quantities one-versus-rest, but we have not tested it. The residual effect of the prevalence at fixed $\nmin$ is small but detectable at weak signal (\cref{sec:exp-ensemble}), and we have no quantitative account of it. Behaviour that depends on library defaults, such as XGBoost's intercept, is specific to the versions we name.

\section{Conclusion}\label{sec:conclusion}

The shadow rule asks whether a feature does better than the best that noise can do. With a rare positive class, how much evidence a feature can show is set by the number of positives and its effect size, through $s = \nmin\delta^2$ or its rank analogue $\sauc$, while the best that noise can do depends on $\nmin$ alone and grows only slowly. The rule's false positive level is unaffected by imbalance; its power is what imbalance removes. For the first tree this is proved, and for full ensembles the experiments find failure rates at least as high as the root-split theory predicts. Two practical consequences follow. Much of what looks like an imbalance effect in XGBoost~1.x and 2.x is the default intercept, which XGBoost~3 has removed and which one argument fixes in earlier versions. And when positives are few, a rank test on each feature uses the same information more efficiently than a tree split and does not pay for a threshold search, at the price of missing signals that only a model can see. Below about 25 positives no method we tested separates weak signal from noise reliably, and the honest answer is to keep features on domain grounds or to collect more positives.
